\documentclass{jsarticle}
\usepackage{amsmath}
\usepackage{amssymb}
\begin{document}
\title{ 日記}
\author{}
\date{}
\maketitle

慣性系は、観測系が固定されている。
非慣性系は、M=1,2,xと接空間が固定されていない。接空間が無数に観測系としてある。
一般共変形は、非共変系から導かれる。M理論をこの非共変系は、エドワード・ウィッテンさんは調べている。アインシュタイン博士は、一時期この非共変系を調べていて、ある時期に一般共変系で一般相対性理論を作り上げた。慣性系と非慣性系を統合したのを一般共変形のことをいう。非共変系は、D-braneとAnti-D-braneを無限に接空間の切片を加群したassemble-D-braneのことをいう。一般共変系は、この非共変系と慣性系を統合したのをいう。異次元空間の一般相対性理論がAdS5の５次元多様体のホログラム空間を表している。深リーマン予想が量子群で、リーマン予想がサーストン多様体のポワンカレ予想のことをいう。ゼータ関数はこのポワンカレ予想のことをいう。
一般相対性理論の場の理論を広中平祐博士はヒルベルト空間の標数０の
体の上の多様体を、大域的微分方程式として、調べていた。
ヒッグス場の方程式は、微分幾何のオイラー数が、なぜ特異点として、
トポロジーとして存在していないかをこのヒルベルト空間のグラスマン多様体を
スピンネットワークとして、ダークマターとビッグバンがなぜ特異点理論として
調べられているかを語っている。このオイラー数と特異点の数値は、
場合分けから、別の性質を物語っている。佐藤先生は、超関数で
この大域的微分方程式を正規部分群として、積分多様体から、
概均質ベクトル空間として、ゼータ関数の性質を調べていた。
微細構造定数と逆関数の関係で、プランク定数とD-braneと
Anti-D-braneの数値の原子間の定数が、AdS5の５次元多様体の方程式として
リサ・ランドール博士がラマン・サンドラム博士と共同研究していた。
ゼータ関数は、この領域で開集合として、レオナルド・オイラーが調べていた。
オイラーの定数は、この量子群とゼータ関数の場の理論の土台になっている。
ニュートンリングは、D-braneとAnti-D-braneの虚数軸で
回転していることを言っている。このD-braneとAnti-D-braneの間に
粒子が流れると虚数軸でD-brane同士が回転する。この粒子がグラビトンらしい。
崩壊定理とは、この非対称性理論を言う。熱力学の第2法則から、
情報エントロピーが増大することを、第1法則からエントロピー不変量を、
この崩壊定理は述べている。ダークマターは、この熱拡散と密度和の機構から、
宇宙の連結和を表している。種数２と特異点の０は違う方程式の性質を言っている。
この数値は、ゼータ関数の場の理論の統合した
ヒッグス場の方程式の部分群の性質を表している。微分幾何の量子化は、
プランク長体積以下になったら、ブラックホールとホワイトホールの
５次元多様体から、スピンネットワークとして、D-braneからもれて
assemble-D-braneをgravitonのcommunicateとして、この次元を伝わっている。
この機構で4種類の力学がDブレーンからもれなかったら、
このボソン力はゼータ関数となる。8種類の微分幾何が、
組み合わせ多様体から、2種類ずつの力の、左巻きと右巻きの遷移元素と同じく、
光量子仮説として、層の理論からこのD-braneが構築されている。
極限値の収束半径の母関数は、加群としての極限値の定数としてあるが、
極限値の加群同士では、収束半径は1種類の定数であるが、
岡本和夫先生は、2種類の力としてのD-braneとAnti-D-braneの加群の
収束値の極限をε・δ論法で、富山大学の数学の教授から教えてもらえた、
$ C^{\infty} $微分可能の積分方程式、
$$ {d \over dx}\int^x_a f(x)dx = f(a) $$
$$ \lim_{x\to a}{f(x)} = l, \lim_{y\to b}
{f(y)} = m, \lim_{x \to a}f(x) + \lim_{y \to b}f(y) = l + m $$

普通は、

$$ \lim_{x,y\to a,b}\{{f(x) + (y)}\} = f(z) $$

となるが、前述の加群とD-braneの方程式は表される。この記述の前にも、
オイラーの公式が、素粒子方程式の場の変換則としてのΓ関数とβ関数の
マトリックスで導かれる、行列式で述べようとして、本書では許容範囲を
超えているとして、述べなかったのが、好いていた。
$$ {d \over df}F = {d \over df}\int\int_M{1 \over (x \log x)^2}dx_m 
+ {d \over df}\int\int_M{1 \over (y \log y)^{1 \over 2}}dy_m $$
$$ e^{i \theta} = \cos \theta + i \sin \theta $$
$$ x^{{1 \over 2} + iy} = e^{x \log x} $$

YMCAの選抜テストでも、このオイラーの公式が出ていた。小川先生と組み合わせ論
のことを相談していたら、横で聞いていた松平先生が、
後から、山口君は理学部に向いていると言ってくれていた。
細木数子からは、丁の己の未と書かれていた。
柳先生の電気技術者の話と理詰めで問題を解く教え方で、YMCAを選んでいた。
代ゼミも本館とサテライトを行き来して、緑のコートで歩き回っていた。
このときは、昼ご飯はちゃんと食べていた。中学2年生と中学3年生のときは、
漫画本を買うために、昼ご飯のお金を貯めていた。このために食べていなかった。
広島大学附属病院東1病棟に入院したときに、よくこれだけの漫画本を買っていたと、
親からあきられていたが、後から、サザエさんの漫画本を親たちはくれていた。
今は、ありったけの漫画本を父親に買ってもらった。
あのときは、よく学生の時に買っていたと思った。あのときは、
サザン・アイズのために歩いていた。代ゼミとYMCAと中学生のときの河合塾は、
先生達から勉強を教えてもらえることに楽しすぎて、歩き回っていた。
すべての勉強は、ハレルヤ学院から始まっていた。
予習と復習、塾と学校、この2校の組み合わせで、勉強が身につく。
オイラーの公式は、人工知能が出来ると言われている。
松平先生からは、英語の成績がぐんぐん上がっていたが、
数学はどんどん成績が下がっていたと言われてた。YMCAの先生達は、
ちゃんと教えていて、戊の寅のときに成績が理系の分野がぐんぐん上がっていた。
私立と2次試験のときの前に上がっていた。私立の東邦大学は、
化学の試験のときに、脱水症状で解き方の式だけ書いて、バタンキューしていた。
全部は伊沢君のおかげで成績がぐんぐん上がっていた。
このときのYMCAとZ-kaiのおかげで2001年の難関国公立大学受験クラスの河合塾と
2006年の京大クラス、SRSも入っている。Photo Readingが出来てよかった。
MindMapは書き出すことができるツールになっていた。
速書法もやっていてよかった。彩さんの本も一冊目の本を読んでいてよかった。
全部で書き出すことが出来るようになった。中学1年生の3学期と高等学校2年生の
3学期のときの姉たちの勉強のアドバイスも入っている。
商代数ができると学問ができるとおもう。順子姉さんの高等学校のときの
私に言いたかったことが、ガロワの姉の好意への思いで知ったことで
すごくわかる。本当にあのときはごめんなさいとしか言えません。
私の大域的微分方程式は、ガロワとライプニッツの式で見つけていた。
お父さんと先生達のヒントも入っている。先生の添削で特異点と
3次元多様体のオイラー数の関係から大域的微分方程式が見つかっている。
グリーシャさんも入っている。
私の運命星は44で運命数は45、宿命数は4、宿命大殺界は34から、始まりは14から、
宿命星が木星の丑だかららしい。金星の丑も入っている。
細木数子は、私を丁の己、と言ってくれていた。
もてないので水星の亥と戌ではないと、宿命星からみんなが言っていた。
手相からいうと土星と木星の霊合らしい。
全部の占いはZappのみんなのおかげです。火星の霊合も入っているらしい。
五行の全部の手相らしい。
お母さんだけが私のことを信頼していた。
細木数子は、アインシュタインかどうか、8月に電話してきて、
駅の子はでたらめと私の誕生日を4月１０日に剣友会の名簿にでたらめを書いていた。
岡先生がそのままにしなさいと言っていたのでそのままにしたら、
みんなはあとあとに私の誕生日を知ったらしく、
なぜ剣道を入会したかを理解していた。舛添要一さんがやめられて、
森さんまでもやめられて、本当に日本の政治家は国民に大切にされているのか
とやりきれないです。謝罪したら、許すのが義務でしょう。
マシーンではないのですよ。誰だってミスをするのですよ。
先生達は私が５０歳までこんなことをできないと言っていたので、今までの人達が
やめさせられたのを今理解しました。
アインシュタインは、蔵王病院から2回目の退院した後に、
順子姉さんと宇多田ヒカルさんで知った。本当は退院ではないらしい。
真面目にアカシックレコードを身につけたら、ゼータ関数の数式が出てきた。
小方君と加藤十吉先生がこのゼータ関数を出した要因に入っている。
全部は岡本和夫先生の集合論のアーベル多様体をゼータ関数で
積分方程式として、導いているのを参考文献にしていたら、
このアカシックレコードにナッシュさんとトビーケリーさんの会話の
非可換代数方程式がきっかけで出てきた。安岡先生は、本当に私の子どもに
名前をつけるのかどうか。私は大変な人生を生きてきたので、
どうかともおもう。頼みたいが頼めるのかどうか。
全部の御霊を入れれるのかどうか。黒には全部入るらしい。
御影は、弘法大師様だけとおもう。YMCAのときに、英語の音読をしていたら、
一回だけ字が、弘法大師様になっていた。伊沢君がこのときのノートを
コピーしてくれた。わたしには無理とおもう。
墓のことは、細木数子にかなう人は安岡先生だけとおもう。
わたしが御先祖様の墓を拭いていたら、家族が修験者みたいと褒めてくれていた。
エロさはどうもとれないとおもう。御先祖様の加護はあるとおもうが。
無理と思う。出来ないと思う。習慣は口癖からと、アフォメーションから、
実現可能らしい。Photo Readingは、最先端の心理学で苫米地英人博士を
先取りしようとしていたが、博士は15歳で作ったらしい。

\end{document}
